
As parameter-efficient fine-tuning methods proliferate, practitioners accumulate collections of LoRA adapters without systematic methods for understanding relationships between them. When selecting an adapter for a new task, current practice relies on metadata tags and manual curation rather than analysis of the adapters themselves. This work investigates whether LoRA adapter weight spaces contain geometric structure that reflects task similarity.

LoRA (Low-Rank Adaptation) constrains weight updates to a low-rank subspace via the decomposition $\Delta W$ = BA, where B $\in \mathbb{R}^(d\times r)$ and A $\in \mathbb{R}^(r\times k)$. The column space of the B matrix defines the output transformation subspace. If semantically similar tasks require similar functional transformations, their B matrices may span similar subspaces. This hypothesis has not been rigorously validated under controlled experimental conditions.

Previous attempts to analyze LoRA adapter geometry used public adapter collections with uncontrolled provenance---different base model checkpoints, quantization levels, LoRA configurations, and training hyperparameters. These uncontrolled variables confound any geometric signal with experimental artifacts. One prior experiment using 8 public HuggingFace adapters observed a large effect size (Cohen's d = 0.91) but insufficient statistical power (p = 0.127) due to the small sample size.

This work addresses these limitations by training adapters in-house under controlled conditions. A single base model checkpoint is fine-tuned with identical LoRA configuration and training hyperparameters across all tasks. The only variable that changes is the training task. This approach, adapted from the Model Zoo methodology for full neural network populations, isolates task-specific geometric structure from experimental confounds.

The contributions of this work are:

\begin{enumerate}
\item An existence proof demonstrating that within-category LoRA adapters exhibit smaller Grassmann distances than between-category adapters under controlled conditions (Cohen's d = 0.765, p = 8.63 $\times 10^{-28})$.
\end{enumerate}

\begin{enumerate}
\item Validation that geometric similarity correlates with semantic similarity as defined by FLAN task taxonomy (Spearman $\rho$ = 0.389, p = 1.29 $\times 10^{-29})$.
\end{enumerate}

\begin{enumerate}
\item A negative finding that task-category clustering is distributed uniformly across all transformer layer types, with no layer type exceeding Cohen's d = 0.8.
\end{enumerate}

\begin{enumerate}
\item Identification of a limitation: the P3 control analysis shows within-task variance ratio = 0.89, indicating that training stochasticity substantially affects adapter geometry.
\end{enumerate}

This work establishes that geometric structure exists in LoRA adapter weight space but does not claim that Grassmann distance outperforms existing task similarity methods. Comparison to baseline methods such as task embedding similarity is deferred to future work.
